\documentclass[10pt,a4paper]{amsart}
\usepackage{amsmath,amssymb,amsthm,graphicx}
\usepackage{verbatim}
\newenvironment{reference}{\comment}{\endcomment}
%\newenvironment{reference}{}{}
\newenvironment{slogan}{\comment}{\endcomment}
\newenvironment{history}{\comment}{\endcomment}

% For commutative diagrams we use Xy-pic
\usepackage[all]{xy}

% We use 2cell for 2-commutative diagrams.
\xyoption{2cell}
% Theorem environments.
%
\theoremstyle{plain}
\newtheorem{theorem}[subsection]{Theorem}
\newtheorem{proposition}[subsection]{Proposition}
\newtheorem{lemma}[subsection]{Lemma}

\theoremstyle{definition}
\newtheorem{definition}[subsection]{Definition}
\newtheorem{example}[subsection]{Example}
\newtheorem{exercise}[subsection]{Exercise}
\newtheorem{situation}[subsection]{Situation}

\theoremstyle{remark}
\newtheorem{remark}[subsection]{Remark}
\newtheorem{remarks}[subsection]{Remarks}

\numberwithin{equation}{subsection}

% Macros
%
\def\lim{\mathop{\mathrm{lim}}\nolimits}
\def\colim{\mathop{\mathrm{colim}}\nolimits}
\def\Spec{\mathop{\mathrm{Spec}}}
\def\Hom{\mathop{\mathrm{Hom}}\nolimits}
\def\Ext{\mathop{\mathrm{Ext}}\nolimits}
\def\SheafHom{\mathop{\mathcal{H}\!\mathit{om}}\nolimits}
\def\SheafExt{\mathop{\mathcal{E}\!\mathit{xt}}\nolimits}
\def\Sch{\mathit{Sch}}
\def\Mor{\mathop{\mathrm{Mor}}\nolimits}
\def\Ob{\mathop{\mathrm{Ob}}\nolimits}
\def\Sh{\mathop{\mathit{Sh}}\nolimits}
\def\NL{\mathop{N\!L}\nolimits}
\def\CH{\mathop{\mathrm{CH}}\nolimits}
\def\proetale{{pro\text{-}\acute{e}tale}}
\def\etale{{\acute{e}tale}}
\def\QCoh{\mathit{QCoh}}
\def\Ker{\mathop{\mathrm{Ker}}}
\def\Im{\mathop{\mathrm{Im}}}
\def\Coker{\mathop{\mathrm{Coker}}}
\def\Coim{\mathop{\mathrm{Coim}}}

% Boxtimes
%
\DeclareMathSymbol{\boxtimes}{\mathbin}{AMSa}{"02}

%
% Macros for moduli stacks/spaces
%
\def\QCohstack{\mathcal{QC}\!\mathit{oh}}
\def\Cohstack{\mathcal{C}\!\mathit{oh}}
\def\Spacesstack{\mathcal{S}\!\mathit{paces}}
\def\Quotfunctor{\mathrm{Quot}}
\def\Hilbfunctor{\mathrm{Hilb}}
\def\Curvesstack{\mathcal{C}\!\mathit{urves}}
\def\Polarizedstack{\mathcal{P}\!\mathit{olarized}}
\def\Complexesstack{\mathcal{C}\!\mathit{omplexes}}
% \Pic is the operator that assigns to X its picard group, usage \Pic(X)
% \Picardstack_{X/B} denotes the Picard stack of X over B
% \Picardfunctor_{X/B} denotes the Picard functor of X over B
\def\Pic{\mathop{\mathrm{Pic}}\nolimits}
\def\Picardstack{\mathcal{P}\!\mathit{ic}}
\def\Picardfunctor{\mathrm{Pic}}
\def\Deformationcategory{\mathcal{D}\!\mathit{ef}}
\usepackage{cleveref}
\crefname{originallemma}{original lemma}{original lemmas}
\Crefname{originallemma}{Original lemma}{Original lemmas}
\makeatletter
\newlabel{constructions-lemma-projective-space}{{01NE}{source}{}{sourcepointer-1}{}}
\newlabel{constructions-lemma-projective-space@cref}{{[originallemma][0][]01NE}{[1][0][]source}{}{}{}}
\makeatother
\crefname{originallemma}{original lemma}{original lemmas}
\Crefname{originallemma}{Original lemma}{Original lemmas}
\makeatletter
\newlabel{constructions-lemma-surjective-graded-rings-generated-degree-1-map-proj}{{01N1}{source}{}{sourcepointer-2}{}}
\newlabel{constructions-lemma-surjective-graded-rings-generated-degree-1-map-proj@cref}{{[originallemma][0][]01N1}{[1][0][]source}{}{}{}}
\makeatother
\crefname{originallemma}{original lemma}{original lemmas}
\Crefname{originallemma}{Original lemma}{Original lemmas}
\makeatletter
\newlabel{divisors-lemma-characterize-OD}{{01X0}{source}{}{sourcepointer-3}{}}
\newlabel{divisors-lemma-characterize-OD@cref}{{[originallemma][0][]01X0}{[1][0][]source}{}{}{}}
\makeatother
\crefname{originallemma}{original lemma}{original lemmas}
\Crefname{originallemma}{Original lemma}{Original lemmas}
\makeatletter
\newlabel{divisors-lemma-finite-trivialize-invertible-upstairs}{{0BUT}{source}{}{sourcepointer-4}{}}
\newlabel{divisors-lemma-finite-trivialize-invertible-upstairs@cref}{{[originallemma][0][]0BUT}{[1][0][]source}{}{}{}}
\makeatother
\crefname{originalremark}{original remark}{original remarks}
\Crefname{originalremark}{Original remark}{Original remarks}
\makeatletter
\newlabel{divisors-remark-ses-regular-section}{{0C6K}{source}{}{sourcepointer-5}{}}
\newlabel{divisors-remark-ses-regular-section@cref}{{[originalremark][0][]0C6K}{[1][0][]source}{}{}{}}
\makeatother
\crefname{originallemma}{original lemma}{original lemmas}
\Crefname{originallemma}{Original lemma}{Original lemmas}
\makeatletter
\newlabel{morphisms-lemma-proper-ample-is-proj}{{0C6J}{source}{}{sourcepointer-6}{}}
\newlabel{morphisms-lemma-proper-ample-is-proj@cref}{{[originallemma][0][]0C6J}{[1][0][]source}{}{}{}}
\makeatother
\crefname{originallemma}{original lemma}{original lemmas}
\Crefname{originallemma}{Original lemma}{Original lemmas}
\makeatletter
\newlabel{varieties-lemma-ample-curve}{{0B5X}{source}{}{sourcepointer-7}{}}
\newlabel{varieties-lemma-ample-curve@cref}{{[originallemma][0][]0B5X}{[1][0][]source}{}{}{}}
\makeatother
\crefname{originallemma}{original lemma}{original lemmas}
\Crefname{originallemma}{Original lemma}{Original lemmas}
\makeatletter
\newlabel{varieties-lemma-chi-tensor-finite}{{0AYT}{source}{}{sourcepointer-8}{}}
\newlabel{varieties-lemma-chi-tensor-finite@cref}{{[originallemma][0][]0AYT}{[1][0][]source}{}{}{}}
\makeatother
\title{On a genus 0 proper curve an invertible sheaf of positive degree d has d+1 sections, vanishing H\textasciicircum{}1, and embeds X into P\textasciicircum{}d}
\author{The Stacks Project Authors}
\date{stacks-project:curves.tex:local-snapshot-20261012; GNU FDL 1.2 or later}
\begin{document}
\maketitle
\noindent\textit{Editorial note.} All delivered mathematical text is the author's own, copied from the pinned Stacks Project chapter file curves.tex; no statement or proof step was rewritten, reordered or removed.\par\medskip
\noindent\textit{Editorial note.} The counted proof is the author's complete argument: the section sequence with its divisor, the trivialization argument on the Artinian zero scheme, the vanishing of H\textasciicircum{}1, the multiplication-map surjectivity with its two-part argument, and the Proj comparison closing the embedding are all retained verbatim, including the author's footnote and his remark that one verification is omitted.\par\medskip
\noindent\textit{Editorial note.} The label lemma-genus-zero-positive-degree names the counted statement, and the proof environment immediately following it in the source is delivered verbatim as its proof.\par\medskip
\noindent\textit{Editorial note.} The counted statement is carried by the author's own lemma environment in the chapter source.\par\medskip
\noindent\textit{Allowed standard prerequisites (editorial).}\begin{itemize}
\item Effective Cartier divisors from sections. Divisors, Lemma divisors-lemma-characterize-OD and Remark divisors-remark-ses-regular-section. Supplies the short exact sequence 0 -> O\_X -> L -> L|\_D -> 0 attached to a nonzero section s with zero scheme D.
\item Ampleness and Proj comparison on a curve. Varieties, Lemma varieties-lemma-ample-curve; Morphisms, Lemma morphisms-lemma-proper-ample-is-proj; Constructions, Lemma constructions-lemma-surjective-graded-rings-generated-degree-1-map-proj. Identifies X with Proj of the section ring and converts the surjectivity of the multiplication maps into the closed immersion into P\textasciicircum{}d\_k.
\end{itemize}\medskip
\ifdefined\hypertarget\hypertarget{sourcepointer-1}{}\fi
\noindent\textit{Source reference.} Original lemma 01NE: sources/stacks-chow/tags.txt, source lines 2155--2155. Prerequisite pointer invoked at proof line(s) 2103; it names a standard fact of a sibling chapter that the counted argument uses, and every step of the counted proof built on it is retained verbatim in the delivered body.\par\medskip
\ifdefined\hypertarget\hypertarget{sourcepointer-2}{}\fi
\noindent\textit{Source reference.} Original lemma 01N1: sources/stacks-chow/tags.txt, source lines 2142--2142. Prerequisite pointer invoked at proof line(s) 2100; it names a standard fact of a sibling chapter that the counted argument uses, and every step of the counted proof built on it is retained verbatim in the delivered body.\par\medskip
\ifdefined\hypertarget\hypertarget{sourcepointer-3}{}\fi
\noindent\textit{Source reference.} Original lemma 01X0: sources/stacks-chow/tags.txt, source lines 2501--2501. Prerequisite pointer invoked at proof line(s) 2042; it names a standard fact of a sibling chapter that the counted argument uses, and every step of the counted proof built on it is retained verbatim in the delivered body.\par\medskip
\ifdefined\hypertarget\hypertarget{sourcepointer-4}{}\fi
\noindent\textit{Source reference.} Original lemma 0BUT: sources/stacks-chow/tags.txt, source lines 14886--14886. Prerequisite pointer invoked at proof line(s) 2054; it names a standard fact of a sibling chapter that the counted argument uses, and every step of the counted proof built on it is retained verbatim in the delivered body.\par\medskip
\ifdefined\hypertarget\hypertarget{sourcepointer-5}{}\fi
\noindent\textit{Source reference.} Original remark 0C6K: sources/stacks-chow/tags.txt, source lines 15300--15300. Prerequisite pointer invoked at proof line(s) 2043; it names a standard fact of a sibling chapter that the counted argument uses, and every step of the counted proof built on it is retained verbatim in the delivered body.\par\medskip
\ifdefined\hypertarget\hypertarget{sourcepointer-6}{}\fi
\noindent\textit{Source reference.} Original lemma 0C6J: sources/stacks-chow/tags.txt, source lines 15299--15299. Prerequisite pointer invoked at proof line(s) 2086; it names a standard fact of a sibling chapter that the counted argument uses, and every step of the counted proof built on it is retained verbatim in the delivered body.\par\medskip
\ifdefined\hypertarget\hypertarget{sourcepointer-7}{}\fi
\noindent\textit{Source reference.} Original lemma 0B5X: sources/stacks-chow/tags.txt, source lines 14050--14050. Prerequisite pointer invoked at proof line(s) 2084; it names a standard fact of a sibling chapter that the counted argument uses, and every step of the counted proof built on it is retained verbatim in the delivered body.\par\medskip
\ifdefined\hypertarget\hypertarget{sourcepointer-8}{}\fi
\noindent\textit{Source reference.} Original lemma 0AYT: sources/stacks-chow/tags.txt, source lines 13801--13801. Prerequisite pointer invoked at proof line(s) 2048; it names a standard fact of a sibling chapter that the counted argument uses, and every step of the counted proof built on it is retained verbatim in the delivered body.\par\medskip
\begin{lemma}
\label{lemma-genus-zero-positive-degree}
Let $X$ be a proper curve over a field $k$ with $H^0(X, \mathcal{O}_X) = k$.
Assume $X$ has genus $0$. Let $\mathcal{L}$ be an invertible
$\mathcal{O}_X$-module of degree $d > 0$. Then we have
\begin{enumerate}
\item $\dim_k H^0(X, \mathcal{L}) = d + 1$ and $\dim_k H^1(X, \mathcal{L}) = 0$,
\item $\mathcal{L}$ is very ample and defines a closed immersion into
$\mathbf{P}^d_k$.
\end{enumerate}
\end{lemma}

\begin{proof}
By definition of degree and genus we have
$$
\dim_k H^0(X, \mathcal{L}) - \dim_k H^1(X, \mathcal{L}) = d + 1
$$
Let $s$ be a nonzero section of $\mathcal{L}$.
Then the zero scheme of $s$ is an effective Cartier
divisor $D \subset X$, we have $\mathcal{L} = \mathcal{O}_X(D)$ and
we have a short exact sequence
$$
0 \to \mathcal{O}_X \to \mathcal{L} \to \mathcal{L}|_D \to 0
$$
see Divisors, Lemma \ref{divisors-lemma-characterize-OD} and
Remark \ref{divisors-remark-ses-regular-section}.
Since $H^1(X, \mathcal{O}_X) = 0$ by assumption, we see that
$H^0(X, \mathcal{L}) \to H^0(X, \mathcal{L}|_D)$ is surjective.
As $\mathcal{L}|_D$ is generated by global sections
(because $\dim(D) = 0$, see
Varieties, Lemma \ref{varieties-lemma-chi-tensor-finite})
we conclude that the invertible module $\mathcal{L}$
is generated by global sections.
In fact, since $D$ is an Artinian scheme we have
$\mathcal{L}|_D \cong \mathcal{O}_D$\footnote{In our case this
follows from Divisors, Lemma
\ref{divisors-lemma-finite-trivialize-invertible-upstairs}
as $D \to \Spec(k)$ is finite.} and hence we can
find a section $t$ of $\mathcal{L}$ whose restriction
of $D$ generates $\mathcal{L}|_D$.
The short exact sequence also shows that $H^1(X, \mathcal{L}) = 0$.

\medskip\noindent
For $n \geq 1$ consider the multiplication map
$$
\mu_n :
H^0(X, \mathcal{L}) \otimes_k H^0(X, \mathcal{L}^{\otimes n})
\longrightarrow
H^0(X, \mathcal{L}^{\otimes n + 1})
$$
We claim this is surjective. To see this we consider the short exact
sequence
$$
0 \to \mathcal{L}^{\otimes n} \xrightarrow{s}
\mathcal{L}^{\otimes n + 1} \to \mathcal{L}^{\otimes n + 1}|_D \to 0
$$
The sections of $\mathcal{L}^{\otimes n + 1}$ coming from the left in this
sequence are in the image of $\mu_n$. On the other hand, since
$H^0(\mathcal{L}) \to H^0(\mathcal{L}|_D)$ is surjective and since
$t^n$ maps to a trivialization of $\mathcal{L}^{\otimes n}|_D$
we see that $\mu_n(H^0(X, \mathcal{L}) \otimes t^n)$ gives a subspace
of $H^0(X, \mathcal{L}^{\otimes n + 1})$ surjecting onto the global sections of
$\mathcal{L}^{\otimes n + 1}|_D$. This proves the claim.

\medskip\noindent
Observe that $\mathcal{L}$ is ample by
Varieties, Lemma \ref{varieties-lemma-ample-curve}.
Hence
Morphisms, Lemma \ref{morphisms-lemma-proper-ample-is-proj}
gives an isomorphism
$$
X \longrightarrow
\text{Proj}\left(
\bigoplus\nolimits_{n \geq 0} H^0(X, \mathcal{L}^{\otimes n})\right)
$$
Since the maps $\mu_n$ are surjective for all $n \geq 1$ we see that
the graded algebra on the right hand side is a quotient of
the symmetric algebra on $H^0(X, \mathcal{L})$. Choosing a $k$-basis
$s_0, \ldots, s_d$ of $H^0(X, \mathcal{L})$ we see that
it is a quotient of a polynomial algebra in $d + 1$ variables.
Since quotients of graded rings correspond to closed immersions
of $\text{Proj}$ (Constructions, Lemma
\ref{constructions-lemma-surjective-graded-rings-generated-degree-1-map-proj})
we find a closed immersion $X \to \mathbf{P}^d_k$. We omit the
verification that this morphism is the morphism of
Constructions, Lemma \ref{constructions-lemma-projective-space}
associated to the sections $s_0, \ldots, s_d$ of $\mathcal{L}$.
\end{proof}

\end{document}
